
\section{$A_\infty$ 代数中的 Whitehead 定理}
\label{sec:4-5}
\liftIndex{I169}

本节证明定理 4.2.45。证明策略与定理 4.2.34 相似：我们将同时
对数目滤过与能量滤过作归纳，构造所需的同伦逆。为使论证清楚，先建立
一种障碍理论。这样，证明就与代数拓扑中熟知的 Whitehead 定理的证明
十分相似。

\subsection{将 $A_K$ 同态延拓为 $A_{K+1}$ 同态}
\label{subsec:4-5-1}

本小节的主要结果是下面的定理 4.5.1。令 $\overline C_i$ 为
$A_{K+1}$ 代数，令
$\overline{\mathfrak f}\colon\overline C_1\to\overline C_2$
为 $A_K$ 同态。考虑 $R$-模
$\operatorname{Hom}(B_{K+1}\overline C_1[1],\overline C_2[1])$。
在此模上定义上边界算子
\[
 \delta_1(\varphi)
 =\overline{\mathfrak m}_1\circ\varphi
 +(-1)^{\deg\varphi+1}
 \varphi\circ\widehat{\overline{\mathfrak m}}_1,
\]
其中
$\widehat{\overline{\mathfrak m}}_1\colon
B\overline C_1[1]\to B\overline C_1[1]$
是由 $\overline C_1$ 上的 $\overline{\mathfrak m}_1$
诱导的余导子，见 (3.2.1)。记
\[
 H\bigl(\operatorname{Hom}
 (B_{K+1}\overline C_1[1],\overline C_2[1]),\delta_1\bigr)
\]
为 $\delta_1$-上同调类的集合。

\begin{theorem}{4.5.1}\label{thm:4-5-1}
存在次数为 $1$ 的
$o_{K+1}(\overline{\mathfrak f})\in
\operatorname{Hom}(B_{K+1}\overline C_1[1],\overline C_2[1])$，
满足以下性质：
\begin{enumerate}[label=\textup{(4.5.2.\arabic*)},leftmargin=5.4em]
 \item $\delta_1(o_{K+1}(\overline{\mathfrak f}))=0$。
 \item 当且仅当存在延拓 $\overline{\mathfrak f}$ 的
  $A_{K+1}$ 同态时\liftIndex{I351}，
  $[o_{K+1}(\overline{\mathfrak f})]=0$ 作为
  $\delta_1$-上同调类。
 \item 若 $\overline{\mathfrak f}'$ 与
  $\overline{\mathfrak f}$ 是 $A_K$ 同伦的，则
  \[
   [o_{K+1}(\overline{\mathfrak f})]
   =[o_{K+1}(\overline{\mathfrak f}')]
   \in H\bigl(\operatorname{Hom}
   (B_{K+1}\overline C_1[1],\overline C_2[1]),\delta_1\bigr).
  \]
 \item 若
  $\overline{\mathfrak g}\colon\overline C_1'\to\overline C_1$、
  $\overline{\mathfrak g}'\colon\overline C_2\to\overline C_2'$
  为 $A_{K+1}$ 同态，则
  \[
   [o_{K+1}(\overline{\mathfrak g}'\circ
      \overline{\mathfrak f}\circ\overline{\mathfrak g})]
   =(\overline{\mathfrak g}'_1)_*
     \circ[o_{K+1}(\overline{\mathfrak f})]
     \circ(\overline{\mathfrak g}_1\otimes\cdots
       \otimes\overline{\mathfrak g}_1)_*
  \]
  在
  $H(\operatorname{Hom}
  (B_{K+1}\overline C_1'[1],\overline C_2'[1]),\delta_1)$
  中成立。其中
  \[
   \overline{\mathfrak g}_1\otimes\cdots
   \otimes\overline{\mathfrak g}_1\colon
   B_{K+1}\overline C_1'[1]\to B_{K+1}\overline C_1[1]
  \]
  由 $\overline{\mathfrak g}_1\colon
  \overline C_1'[1]\to\overline C_1[1]$ 诱导；它又在
  $\delta_1$-上同调上诱导同态
  $(\overline{\mathfrak g}_1\otimes\cdots
  \otimes\overline{\mathfrak g}_1)_*$。
\end{enumerate}
\end{theorem}

\begin{remark}{4.5.3}\label{rem:4-5-3}
事实上，在证明定理 4.2.34 时，我们已用到类似结论。由于定理
4.5.1 的陈述中出现了同伦概念，我们在此重复论证，以表明论证并非
循环论证。
\end{remark}

\begin{proof}[定理 4.5.1 的证明]
先取 $\overline{\mathfrak f}'_{K+1}=0$，考虑
\[
 \widehat{\overline{\mathfrak f}}_{1\cdots K+1}
 \in\operatorname{Hom}
 (B_{1\cdots K+1}\overline C_1[1],
  B_{1\cdots K+1}\overline C_2[1]).
\]
由引理 4.4.9 和 4.4.10，它诱导上链同态
$\widehat{\overline{\mathfrak f}}_{2\cdots K+1}$ 与
$\widehat{\overline{\mathfrak f}}_{1\cdots K}$。因此，
$\widehat{\overline d}_{1\cdots K+1}
\circ\widehat{\overline{\mathfrak f}}_{1\cdots K+1}
-\widehat{\overline{\mathfrak f}}_{1\cdots K+1}
\circ\widehat{\overline d}_{1\cdots K+1}$
在 $B_{1\cdots K}\overline C_1[1]$ 上为零，而其像落在
$B_1\overline C_2[1]=\overline C_2[1]$ 中。定义
\[
 o_{K+1}(\overline{\mathfrak f})
 =\widehat{\overline d}_{1\cdots K+1}
   \circ\widehat{\overline{\mathfrak f}}_{1\cdots K+1}
  -\widehat{\overline{\mathfrak f}}_{1\cdots K+1}
   \circ\widehat{\overline d}_{1\cdots K+1}
 \in\operatorname{Hom}
 (B_{K+1}\overline C_1[1],\overline C_2[1]).
\]
于是
\[
 \delta_1(o_{K+1}(\overline{\mathfrak f}))
 =\overline{\mathfrak m}_1\circ o_{K+1}(\overline{\mathfrak f})
 -o_{K+1}(\overline{\mathfrak f})
  \circ\widehat{\overline{\mathfrak m}}_1=0,
\]
得到 (4.5.2.1)。

要证 (4.5.2.2)，只须注意：
$\overline{\mathfrak f}_i$（$i=1,\ldots,K$）与
$\overline{\mathfrak f}_{K+1}$ 合成一个 $A_{K+1}$ 同态，当且仅当
$\delta_1(\overline{\mathfrak f}_{K+1})
=-o_{K+1}(\overline{\mathfrak f})$。

再证 (4.5.2.4)。考虑由
$\overline{\mathfrak g}_i$（$i=1,\ldots,K$）诱导的余代数同态
\[
 \widehat{\overline{\mathfrak g}}\colon
 B_{1\cdots K+1}\overline C'_1[1]\longrightarrow
 B_{1\cdots K+1}\overline C_1[1],
 \qquad
 \widehat{\overline{\mathfrak g}}'\colon
 B_{1\cdots K+1}\overline C_2[1]\longrightarrow
 B_{1\cdots K+1}\overline C'_2[1].
\]
带帽延拓与复合可交换。现考察交换子
\[
 \{\widehat d,\widehat{\overline{\mathfrak g}}'
    \circ\widehat{\overline{\mathfrak f}}
    \circ\widehat{\overline{\mathfrak g}}\}
 =\{\widehat d,\widehat{\overline{\mathfrak g}}'\}
    \circ\widehat{\overline{\mathfrak f}}
    \circ\widehat{\overline{\mathfrak g}}
 +\widehat{\overline{\mathfrak g}}'
    \circ\{\widehat d,\widehat{\overline{\mathfrak f}}\}
    \circ\widehat{\overline{\mathfrak g}}
 +\widehat{\overline{\mathfrak g}}'
    \circ\widehat{\overline{\mathfrak f}}
    \circ\{\widehat d,\widehat{\overline{\mathfrak g}}\}.
\]
左端的
$\operatorname{Hom}(B_{K+1}\overline C'_1[1],
\overline C'_2[1])$ 分量是
$[o_{K+1}(\widehat{\overline{\mathfrak g}}'
\circ\widehat{\overline{\mathfrak f}}
\circ\widehat{\overline{\mathfrak g}})]$。
右端第一项是
$-\delta_1\overline{\mathfrak g}'_{K+1}
\circ\overline{\mathfrak f}_1^{\otimes(K+1)}
\circ\overline{\mathfrak g}_1^{\otimes(K+1)}$，
因而是 $\delta_1$-上边界；第二项是
$\overline{\mathfrak g}'_1\circ
o_{K+1}(\overline{\mathfrak f})\circ
(\overline{\mathfrak g}_1\otimes\cdots
\otimes\overline{\mathfrak g}_1)$；第三项是
$-\overline{\mathfrak g}'_1\circ
\overline{\mathfrak f}_1\circ
\delta_1\overline{\mathfrak g}_{K+1}$，
也为 $\delta_1$-上边界。这就得到 (4.5.2.4)。

要证 (4.5.2.3)，还需以下引理。
\begin{lemma}{4.5.4}\label{lem:4-5-4}
若 $\overline{\mathfrak f}\colon\overline C_1\to\overline C_2$
与 $\overline{\mathfrak g}\colon\overline C_1\to\overline C_2$
是 $A_1$ 同伦的，则
\[
 \overline{\mathfrak f}_*=\overline{\mathfrak g}_*\colon
 H^*(\overline C_1;\overline{\mathfrak m}_1)
 \longrightarrow
 H^*(\overline C_2;\overline{\mathfrak m}_1).
\]
\end{lemma}
\begin{proof}
令 $\overline{\mathfrak h}\colon\overline C_1\to
\overline{\mathfrak C}_2$ 为
$\overline{\mathfrak f}$ 与 $\overline{\mathfrak g}$
之间的 $A_1$ 同伦。由 (4.2.2.3)，
$\overline{\mathrm{Incl}}$ 在
$\overline{\mathfrak m}_1$-上同调上诱导同构；
另一方面，由 (4.2.2.2)，对任意 $s_0$ 都有
$\overline{\mathrm{Eval}}_{s=s_0}
\circ\overline{\mathrm{Incl}}=\mathrm{id}$。因此，在上同调上
$\overline{\mathrm{Eval}}_{s=0,*}
=\overline{\mathrm{Eval}}_{s=1,*}\colon
H^*(\overline{\mathfrak C}_2;\overline{\mathfrak m}_1)
\to H^*(\overline C_2;\overline{\mathfrak m}_1)$。从而
$\overline{\mathfrak f}_*
=\overline{\mathrm{Eval}}_{s=0,*}
\circ\overline{\mathfrak h}_*
=\overline{\mathrm{Eval}}_{s=1,*}
\circ\overline{\mathfrak h}_*
=\overline{\mathfrak g}_*$，如所需。
\end{proof}

现在证明 (4.5.2.3)。令
$\overline{\mathfrak h}\colon\overline C_1
\to\overline{\mathfrak C}_2$
为 $\overline{\mathfrak f}$ 与
$\overline{\mathfrak f}'$ 之间的 $A_K$ 同伦。
由 (4.5.2.4) 及引理 4.5.4，有
\[
 \begin{aligned}
 [o_{K+1}(\overline{\mathfrak f})]
 &=[o_{K+1}(\overline{\mathrm{Eval}}_{s=0}
       \circ\overline{\mathfrak h})]
  =\overline{\mathrm{Eval}}_{s=0,*}
       [o_{K+1}(\overline{\mathfrak h})]\\
 &=\overline{\mathrm{Eval}}_{s=1,*}
       [o_{K+1}(\overline{\mathfrak h})]
  =[o_{K+1}(\overline{\mathrm{Eval}}_{s=1}
       \circ\overline{\mathfrak h})]
  =[o_{K+1}(\overline{\mathfrak f}')].
 \end{aligned}
\]
定理 4.5.1 得证。
\end{proof}

下一小节要用到以下推论。
\begin{corollary}{4.5.5}\label{cor:4-5-5}
设 $\overline{\mathfrak f}\colon\overline C_1\to\overline C_2$
是 $A_{K+1}$ 同态，
$\overline{\mathfrak g}\colon\overline C_1\to\overline C_2$
是 $A_K$ 同态，并设
$\overline{\mathfrak h}\colon\overline C_1
\to\overline{\mathfrak C}_2$
为从 $\overline{\mathfrak f}$ 到 $\overline{\mathfrak g}$
的 $A_K$ 同伦。则 $\overline{\mathfrak g}$ 可延拓为
$A_{K+1}$ 同态 $\overline{\mathfrak g}'$，而
$\overline{\mathfrak h}$ 可延拓为从
$\overline{\mathfrak f}$ 到 $\overline{\mathfrak g}'$
的 $A_{K+1}$ 同伦。
\end{corollary}

\begin{proof}
由 (4.5.2.2)、(4.5.2.4) 可知
$(\overline{\mathrm{Eval}}_{s=0})_{1*}
[o_{K+1}(\overline{\mathfrak h})]
=[o_{K+1}(\overline{\mathfrak f})]=0$；
而 $(\overline{\mathrm{Eval}}_{s=0})_{1*}$
在上同调上诱导同构，所以
$\overline{\mathfrak h}$ 可延拓为到区间模型的
$A_{K+1}$ 同态。
为证推论，所取延拓还须满足一些附加性质，下面加以构造。

取
$\overline{\mathfrak h}'_{K+1}
=\overline{\mathrm{Incl}}\circ
\overline{\mathfrak f}_{K+1}$。于是
\[
 (\overline{\mathrm{Eval}}_{s=0})_1
 \bigl(o_{K+1}(\overline{\mathfrak h})
       +\delta_1(\overline{\mathfrak h}'_{K+1})\bigr)=0,
 \qquad
 (\overline{\mathrm{Eval}}_{s=0})_1
 \circ\overline{\mathfrak h}'_{K+1}
 =\overline{\mathfrak f}_{K+1}.
\]
由区间模型的 (4.2.2.3)、(4.2.2.4)，相应上链映射
$(\overline{\mathrm{Eval}}_{s=0})_1$
分裂满射且诱导上同调同构，故其核是
$\delta_1$-无环复形。\translatornote{原书此处把核写成
“$\overline{\mathrm{Eval}}_{s=0,*}$ 的核”，但星号在紧前表示
上同调诱导映射；需要无环性的是它在相关 $\operatorname{Hom}$
复形上诱导的上链映射之核。}
因此可取 $\Delta\overline{\mathfrak h}_{K+1}$，使
\[
 o_{K+1}(\overline{\mathfrak h})
 +\delta_1(\overline{\mathfrak h}'_{K+1})
 =\delta_1(\Delta\overline{\mathfrak h}_{K+1}),
 \qquad
 (\overline{\mathrm{Eval}}_{s=0})_1
 (\Delta\overline{\mathfrak h}_{K+1})=0.
\]
置
$\overline{\mathfrak h}_{K+1}
=\overline{\mathfrak h}'_{K+1}
-\Delta\overline{\mathfrak h}_{K+1}$。
于是 $\overline{\mathfrak h}_k$（$k=1,\ldots,K+1$）
给出 $A_{K+1}$ 同态
$\widetilde{\overline{\mathfrak h}}$。
再定义
$\overline{\mathfrak g}'
=(\overline{\mathrm{Eval}}_{s=1})_1
\circ\widetilde{\overline{\mathfrak h}}$。
容易看出，它是延拓
$\overline{\mathfrak g}$ 的 $A_{K+1}$ 同态。
又因
$(\overline{\mathrm{Eval}}_{s=0})_1
\circ\overline{\mathfrak h}_{K+1}
=\overline{\mathfrak f}_{K+1}$，
故 $\widetilde{\overline{\mathfrak h}}$
是从 $\overline{\mathfrak f}$ 到
$\overline{\mathfrak g}'$ 的 $A_{K+1}$ 同伦。
\end{proof}
\subsection{定理 4.2.45 的证明 I：数目滤过}
\label{subsec:4-5-2}
\liftIndex{I094}

设 $\overline{\mathfrak f}\colon\overline C_1\to\overline C_2$
是弱同伦等价。

\begin{proposition}{4.5.6}\label{prop:4-5-6}
设
$\overline{\mathfrak g}^{(K)}\colon
\overline C_2\to\overline C_1$
为 $A_K$ 同态，
$\overline{\mathfrak h}^{(K)}\colon
\overline C_1\to\overline{\mathfrak C}_1$
为从恒等同态到
$\overline{\mathfrak g}^{(K)}
\circ\overline{\mathfrak f}$ 的 $A_K$ 同伦。
则 $\overline{\mathfrak g}^{(K)}$ 可延拓为
$A_{K+1}$ 同态 $\overline{\mathfrak g}^{(K+1)}$，
且 $\overline{\mathfrak h}^{(K)}$ 可延拓为
从恒等同态到
$\overline{\mathfrak g}^{(K+1)}
\circ\overline{\mathfrak f}$ 的 $A_{K+1}$ 同伦
$\overline{\mathfrak h}^{(K+1)}$。
\end{proposition}

\begin{center}
\resizebox{0.96\linewidth}{!}{$
\begin{array}{ccccccc}
 \overline C_1
 &\xleftarrow{\ \overline{\mathrm{Eval}}_{s=0}\ }&
 \overline{\mathfrak C}_1
 &\xrightarrow{\ \overline{\mathrm{Eval}}_{s=1}\ }&
 \overline C_1
 &\xleftarrow{\ \overline{\mathfrak g}^{(K)}\ }&
 \overline C_2\\[0.5em]
 \Vert&&\underset{\scriptstyle\overline{\mathfrak h}}{\Big\uparrow}
 &&&&\Vert\\[0.5em]
 \overline C_1
 &\mathrel{\hspace{1.7em}=\hspace{1.7em}}&
 \overline C_1
 &\mathrel{\hspace{1.7em}=\hspace{1.7em}}&
 \overline C_1
 &\xrightarrow{\ \overline{\mathfrak f}\ }&
 \overline C_2
\end{array}
$}
\par\smallskip\textbf{图式 4.5.1}\label{diag:4-5-1}
\end{center}

\begin{proof}
由推论 4.5.5，可取
$\overline{\mathfrak h}'_{K+1}$，使
$\overline{\mathfrak h}_k$（$k=1,\ldots,K$）与
$\overline{\mathfrak h}'_{K+1}$ 给出从
$\overline C_1$ 到 $\overline{\mathfrak C}_1$
的 $A_{K+1}$ 同态
$\overline{\mathfrak h}^{(K+1)\prime}$，且
$\overline{\mathrm{Eval}}_{s=0}
\circ\overline{\mathfrak h}^{(K+1)\prime}
=\mathrm{id}$。
另一方面，
$\overline{\mathrm{Eval}}_{s=1}
\circ\overline{\mathfrak h}'_{K+1}$
把
$\overline{\mathfrak g}^{(K)}
\circ\overline{\mathfrak f}$
延拓为一个 $A_{K+1}$ 同态。具体说，
\begin{equation}
 o_{K+1}(\overline{\mathfrak g}^{(K)}
          \circ\overline{\mathfrak f})
 =-\delta_1\bigl(
   \overline{\mathrm{Eval}}_{s=1}
   \circ\overline{\mathfrak h}'_{K+1}\bigr).
\tag{4.5.7}\label{eq:4-5-7}
\end{equation}
由于 $\overline{\mathfrak f}\colon\overline C_1\to\overline C_2$
是弱同伦等价，由 (4.5.2.4) 和 \eqref{eq:4-5-7} 可知
$[o_{K+1}(\overline{\mathfrak g}^{(K)})]=0$。
故可取 $\overline{\mathfrak g}'_{K+1}$，使
\begin{equation}
 o_{K+1}(\overline{\mathfrak g}^{(K)})
 =-\delta_1(\overline{\mathfrak g}'_{K+1}).
\tag{4.5.8}\label{eq:4-5-8}
\end{equation}

定义
\[
 \Xi=
 \overline{\mathfrak g}'_{K+1}
 \circ\overline{\mathfrak f}_1^{\otimes(K+1)}
 -\overline{\mathrm{Eval}}_{s=1}
 \circ\overline{\mathfrak h}'_{K+1}
 \in\operatorname{Hom}
 (B_{K+1}\overline C_1[1],\overline C_1[1]).
\]
由 \eqref{eq:4-5-7} 有 $\delta_1(\Xi)=0$。
由于 $\overline{\mathfrak f}$ 是弱同伦等价，存在
$\delta_1$-闭上链
$\Delta\overline{\mathfrak g}'_{K+1}\in
\operatorname{Hom}(B_{K+1}\overline C_2[1],\overline C_1[1])$，
使
$[\Xi+
(\Delta\overline{\mathfrak g}'_{K+1}
\circ\overline{\mathfrak f}_1^{\otimes(K+1)})]=0$。
因此，存在
$\Delta_1\overline{\mathfrak h}_{K+1}
\in\operatorname{Hom}
(B_{K+1}\overline C_1[1],\overline C_1[1])$，使
\begin{equation}
 \delta_1(\Delta_1\overline{\mathfrak h}_{K+1})
 =(\overline{\mathfrak g}'_{K+1}
   +\Delta\overline{\mathfrak g}'_{K+1})
   \circ\overline{\mathfrak f}_1^{\otimes(K+1)}
  -\overline{\mathrm{Eval}}_{s=1}
   \circ\overline{\mathfrak h}'_{K+1}.
\tag{4.5.9}\label{eq:4-5-9}
\end{equation}

由 (4.2.2.4)，可取
$\Delta\overline{\mathfrak h}_{K+1}
\in\operatorname{Hom}
(B_{K+1}\overline C_1[1],\overline{\mathfrak C}_1[1])$，
使
\begin{equation}
 (\overline{\mathrm{Eval}}_{s=0})_1
  \circ\Delta\overline{\mathfrak h}_{K+1}=0,
 \qquad
 (\overline{\mathrm{Eval}}_{s=1})_1
  \circ\Delta\overline{\mathfrak h}_{K+1}
  =\Delta_1\overline{\mathfrak h}_{K+1}.
\tag{4.5.10}\label{eq:4-5-10}
\end{equation}
现置
$\overline{\mathfrak g}_{K+1}
=\overline{\mathfrak g}'_{K+1}
+\Delta\overline{\mathfrak g}'_{K+1}$，
$\overline{\mathfrak h}_{K+1}
=\overline{\mathfrak h}'_{K+1}
+\delta_1(\Delta\overline{\mathfrak h}_{K+1})$。
下面核对所需性质。

由于 $\Delta\overline{\mathfrak g}'_{K+1}$ 是
$\delta_1$-闭上链，由 \eqref{eq:4-5-8} 得
\[
 o_{K+1}(\overline{\mathfrak g}^{(K)})
 =-\delta_1(\overline{\mathfrak g}_{K+1}).
\]
故 $\overline{\mathfrak g}_{K+1}$ 将
$\overline{\mathfrak g}^{(K)}$ 延拓成
$A_{K+1}$ 同态
$\overline{\mathfrak g}^{(K+1)}$。同样，
$\overline{\mathfrak h}_{K+1}$ 将
$\overline{\mathfrak h}^{(K)}$ 延拓为
$A_{K+1}$ 同态
$\overline{\mathfrak h}^{(K+1)}$。
由 \eqref{eq:4-5-10} 有
\[
 (\overline{\mathrm{Eval}}_{s=0})_1
 \circ\overline{\mathfrak h}^{(K+1)}
 =(\overline{\mathrm{Eval}}_{s=0})_1
 \circ\overline{\mathfrak h}^{(K+1)\prime}
 =\mathrm{id}.
\]
最后，由 \eqref{eq:4-5-9} 和
\eqref{eq:4-5-10}，得到
\[
 (\overline{\mathrm{Eval}}_{s=1})_1
 \circ\overline{\mathfrak h}_{K+1}
 =(\overline{\mathrm{Eval}}_{s=1})_1
  \circ\overline{\mathfrak h}'_{K+1}
  +\delta_1(\Delta_1\overline{\mathfrak h}_{K+1})
 =\overline{\mathfrak g}_{K+1}
  \circ\overline{\mathfrak f}_1^{\otimes(K+1)}.
\]
因此，
$\overline{\mathfrak h}^{(K+1)}$
是从 $\mathrm{id}$ 到
$\overline{\mathfrak g}^{(K+1)}
\circ\overline{\mathfrak f}$ 的
$A_{K+1}$ 同伦。
\end{proof}

现在可以完成定理 4.2.45 (1) 的证明。令
$\overline{\mathfrak f}\colon\overline C_1
\to\overline C_2$ 为弱同伦等价。
由于
$\overline{\mathfrak f}_1\colon
(\overline C_1,\overline{\mathfrak m}_1)
\to(\overline C_2,\overline{\mathfrak m}_1)$
是上链同伦等价，存在上链映射
$\overline{\mathfrak g}_1\colon
(\overline C_2,\overline{\mathfrak m}_1)
\to(\overline C_1,\overline{\mathfrak m}_1)$，
使
$\overline{\mathfrak g}_1
\circ\overline{\mathfrak f}_1$
与恒等映射上链同伦。
$\overline{\mathfrak g}_1$
定义 $A_1$ 同态
$\overline{\mathfrak g}^{(1)}\colon
\overline C_2\to\overline C_1$。
令 $\overline{\mathfrak h}^{a}_1\colon
\overline C_1\to\overline C_1$
为从恒等映射到
$\overline{\mathfrak g}_1\circ
\overline{\mathfrak f}_1$
的上链同伦。
利用 (4.2.2.4) 可取
$\overline{\mathfrak h}^{b}_1\colon
\overline C_1\to\overline{\mathfrak C}_1$，使
\[
 (\overline{\mathrm{Eval}}_{s=1})_1
  \circ\overline{\mathfrak h}^{b}_1
  =\overline{\mathfrak h}^{a}_1,
 \qquad
 (\overline{\mathrm{Eval}}_{s=0})_1
  \circ\overline{\mathfrak h}^{b}_1=0.
\]
置
\[
 \overline{\mathfrak h}_1
 =(\overline{\mathrm{Incl}})_1
  +\overline{\mathfrak m}_1\circ
   \overline{\mathfrak h}^{b}_1
  +\overline{\mathfrak h}^{b}_1\circ
   \overline{\mathfrak m}_1
 \colon\overline C_1\to\overline{\mathfrak C}_1.
\]
容易验证，
$\overline{\mathfrak h}_1$
定义从 $\mathrm{id}$ 到
$\overline{\mathfrak g}^{(1)}
\circ\overline{\mathfrak f}$ 的 $A_1$ 同伦。

对命题 4.5.6 归纳应用，得到
$A_\infty$ 同态
$\overline{\mathfrak g}\colon
\overline C_2\to\overline C_1$，
以及从 $\mathrm{id}$ 到
$\overline{\mathfrak g}
\circ\overline{\mathfrak f}$
的 $A_\infty$ 同伦
$\overline{\mathfrak h}$。

最后证明
$\overline{\mathfrak f}
\circ\overline{\mathfrak g}$
也同伦于恒等同态。
注意 $\overline{\mathfrak g}$ 是弱同伦等价，因为
$\overline{\mathfrak g}_1$
与 $\overline{\mathfrak f}_1$ 的上链同伦逆上链同伦。
于是可再次应用上述论证，得到
$A_\infty$ 同态 $\overline{\mathfrak f}'$，
使 $\overline{\mathfrak f}'
\circ\overline{\mathfrak g}$
同伦于恒等同态。由引理 4.2.43，
$\overline{\mathfrak f}'
\circ\overline{\mathfrak g}
\circ\overline{\mathfrak f}$
同伦于 $\overline{\mathfrak f}$；
另一方面，
$\overline{\mathfrak g}
\circ\overline{\mathfrak f}$
同伦于恒等同态，同一引理又表明
$\overline{\mathfrak f}'
\circ\overline{\mathfrak g}
\circ\overline{\mathfrak f}$
同伦于 $\overline{\mathfrak f}'$。
故 $\overline{\mathfrak f}'$
与 $\overline{\mathfrak f}$ 同伦。
再由
$\overline{\mathfrak f}'
\circ\overline{\mathfrak g}$
同伦于恒等同态，便知
$\overline{\mathfrak f}
\circ\overline{\mathfrak g}$
也同伦于恒等同态。
定理 4.2.45 (1) 证毕。
